\section{Low-energy families on prescribed intervals}

\begin{theorem}[The capped width and its logarithm]
    \label{thm:low_energy_cap}
    \lean{Real.low_energy_cap_bound}
    \leanok
    If $w>0$ and $6w\le\eta_{\mathrm{LD}}$, set
    $g=\min\{w,1\}$ and $M=\max\{1,\eta_{\mathrm{LD}}/6\}$.
    Then $g^{-1}\le M/w$ and
    $\log(g^{-1})=\max\{\log(w^{-1}),0\}$.
\end{theorem}
\begin{proof}
    \leanok
    If $w\le1$, then $g=w$ and $M\ge1$.
    If $w>1$, then $g=1$ and $w\le\eta_{\mathrm{LD}}/6\le M$.
\end{proof}

\begin{theorem}[An explicit count at the source scale]
    \label{thm:low_energy_source_count}
    \lean{Real.low_energy_count_le_source_scale}
    \leanok
    Let $L\ge2$, $q\ge1$, $\mu>0$, and $K\ge0$. Set
    $w=\mu/(3\log L)$ and $g=\min\{w,1\}$, and suppose
    $6w\le\eta_{\mathrm{LD}}$. If a count satisfies
    \begin{align*}
        \log D\le\frac K g[1+\log q+\log(1/g)]^3,
    \end{align*}
    then, with $M=\max\{1,\eta_{\mathrm{LD}}/6\}$,
    \begin{align}
        \log D\le 3KM\frac{\log L}{\mu}
        \left[1+\log q+
          \max\left\{\log\left(\frac{3\log L}{\mu}\right),0\right\}\right]^3.
        \label{eq:ld_source_scale_count}
    \end{align}
\end{theorem}
\begin{proof}
    \leanok
    \uses{thm:low_energy_cap}
    Apply the two preceding bounds and use $1/w=3\log L/\mu$.
\end{proof}

\begin{theorem}[The joint scalar count and strict contraction]
    \label{thm:low_energy_joint_scalar_parameters}
    \lean{Real.exists_low_energy_joint_scalar_parameters}
    \leanok
    Under the numerical assumptions of
    Theorem~\ref{thm:low_energy_scalar_parameters}, the constant can be
    chosen so that its logarithmic bound and strict tradeoff hold at the
    larger count $D=(q\ell^2)^{22\ell}$. The lower bound $\ell\ge21$
    and logarithmic estimate hold for every $q\ge2$ and $0<g\le1$ at the
    chosen scale, before any physical threshold or ground energy is given.
\end{theorem}
\begin{proof}
    \leanok
    \uses{thm:low_energy_threshold_bound}
    The denominator has the same norm bound as before. The strict width
    inequality $w\ge g>4g/5$ makes the physical contraction strictly
    smaller than the scalar comparison contraction. Thus the full joint
    count retains the strict product bound.
\end{proof}

\begin{theorem}[One spanning family on a short region]
    \label{thm:short_region_spanning_family}
    \lean{TNLean.MPS.AreaLaw.exists_short_actual_region_spanning_family}
    \leanok
    Let $q\ge2$, $\ell\ge1$ and $|J|\le4\ell$, and put
    $D=(q\ell^2)^{22\ell}$. For every finite complementary Hilbert
    space, there is a family $(F_j)_{j=1}^{D^2}$ on $\mathcal H_J$,
    chosen before every target subspace. If a nonzero target has a
    $\delta$-viable subspace $S\subseteq\mathcal H_J$ with $\delta<1$,
    then $\sum_jF_jS=\mathcal H_J$, and this enlargement is zero-viable.
\end{theorem}
\begin{proof}
    \leanok
    Configuration matrix units send every nonzero subspace onto a family
    spanning the entire regional space. Viability below one for a nonzero
    target forces $S\ne0$. There are $q^{2|J|}\le q^{8\ell}$ matrix
    units; extend their family by zero to the stated count.
\end{proof}

\begin{theorem}[Enlarging the five-region target region]
    \label{thm:low_energy_regional_family}
    \lean{TNLean.MPS.AreaLaw.exists_five_region_low_energy_family_on_larger_region}
    \leanok
    Let $q\ge2$, $\ell\ge1$, $a+2\ell\le b$, and
    $b+2\ell\le L-1$. Let $\varepsilon,w>0$ and $3w<\eta$, with
    supported positive unit nearest-neighbour terms as in
    Theorem~\ref{thm:physical_low_energy_filter_family}. Put
    $J_0=\{a+\ell,\ldots,b+\ell\}$. Let $J$ be any set of sites with
    $J_0\subseteq J$ and $|J\setminus J_0|\le2\ell$, and set
    $D=(q\ell^2)^{22\ell}$. The attained regional energies, threshold,
    centered original Hamiltonian and clipped Hamiltonian can be chosen
    with the same bounds and both spectral comparisons as in that theorem.
    On $\mathcal H_J$ there is one family $(W_j)_{j=1}^N$ with $N\le D^2$,
    chosen before every viable subspace. If $S$ is $\delta$-viable for the
    clipped upper target and $\delta<1$, then $\sum_jW_jS$ has dimension
    at most $N\dim S$ and is $\Delta/(1-\delta)^2$-viable for the clipped
    lower target, with the same explicit $\Delta$ as in that theorem.
\end{theorem}
\begin{proof}
    \leanok
    \uses{thm:physical_low_energy_filter_family,
      thm:actual_clipped_polynomial_family,thm:chebyshev_nested_polynomial_family}
    Keep the physical truncation and its two spectral comparisons.
    Decomposing its degree-$\ell^2$ polynomial on $J$ multiplies the
    squared core count by at most $q^{2|J\setminus J_0|}$.
    The collar bound gives
    $(q\ell^2)^{40\ell}q^{4\ell}\le D^2$.
    Apply the common-polynomial viability estimate to this family and
    the same clipped Hamiltonian.
\end{proof}

\begin{theorem}[A low-energy improving family on every interval]
    \label{thm:physical_low_energy_interval_family}
    \lean{TNLean.MPS.AreaLaw.exists_physical_low_energy_interval_family}
    \leanok
    Fix $\eta_{\mathrm{LD}}>0$. A constant $A\ge21$ can be chosen before
    $L\ge2$, $q\ge2$, $\mu>0$, the bandwidth
    $2\mu/\log L\le\eta_1\le\eta_{\mathrm{LD}}$, and the supported
    positive unit nearest-neighbour Hamiltonian. Put
    \begin{align*}
        w&=\mu/(3\log L),& g&=\min\{w,1\},&
        u&=1+\log q+\log(1/g),\\
        \ell&=\lceil Au^2/g\rceil,& D&=(q\ell^2)^{22\ell},&
        K&=22(A+1)[6+2\log(A+1)].
    \end{align*}
    Then $\ell\ge21$, and $D$ satisfies
    (\ref{eq:ld_source_scale_count}).

    For every prescribed interval $J=[p,r)\subseteq[0,L)$, including
    empty intervals and intervals meeting either endpoint, there are a
    centered original Hamiltonian $C=H-s\Id$ with an attained ground
    energy $E_0$, nested subspaces $T_-\subseteq T_+$, a positive
    parameter $\Delta$ and a family of $D^2$ operators on $\mathcal H_J$.
    The original spectral space below $E_0+\eta_1$ is $1/200$-close to
    $T_+$, and $T_-$ is $1/200$-close to the original space below
    $E_0+\eta_1-3w$. Moreover $D^{12}\Delta<10^{-5}$.
    These choices precede every viable subspace. If $S$ is
    $\delta$-viable for $T_+$ with $\delta<1$, then $\sum_jF_jS$ has
    dimension at most $D^2\dim S$ and is
    $\Delta/(1-\delta)^2$-viable for $T_-$.
    This supports \cite[Theorem~4.8]{Arad2016RigorousRG} with the finite
    parameters displayed above.
\end{theorem}
\begin{proof}
    \leanok
    \uses{thm:low_energy_source_count,thm:low_energy_joint_scalar_parameters,
      thm:short_region_spanning_family,thm:low_energy_regional_family}
    For $|J|>4\ell$, inward effective cuts give a five-region core
    $J_0\subseteq J$ with $|J\setminus J_0|\le2\ell$. Apply the actual
    physical truncation construction and decompose its polynomial on $J$.
    The regional-family count gains at most $q^{2|J\setminus J_0|}$,
    which fits the joint count. Its strict width margin gives the stated
    tradeoff. Extending the family by zero preserves every image span.

    For $|J|\le4\ell$, use the existing matrix units and choose
    $\Delta=1/(200000D^{12})$. Take $T_+$ and $T_-$ to be the two
    original spectral spaces. The zero target is included; otherwise
    the enlarged regional space is the whole space and has zero error.
    Thus both cases give the same dimension and viability conclusions.
\end{proof}

\begin{theorem}[The same family on the original spectral windows]
    \label{thm:physical_original_low_energy_interval_family}
    \lean{TNLean.MPS.AreaLaw.exists_physical_original_low_energy_interval_family}
    \leanok
    All choices in Theorem~\ref{thm:physical_low_energy_interval_family}
    can be retained simultaneously with the following additional conclusion.
    If $S$ is $\delta$-viable for the original upper window,
    $\delta\ge0$ and $2(\delta+1/200)<1$, then $\sum_jF_jS$ is viable
    for the original lower window with error
    \begin{align*}
        2\left(\frac{\Delta}{[1-2(\delta+1/200)]^2}+\frac1{200}\right).
    \end{align*}
    Its dimension is still at most $D^2\dim S$. The improving-family
    parameter and the final original-window error remain distinct.
\end{theorem}
\begin{proof}
    \leanok
    \uses{thm:physical_low_energy_interval_family}
    Apply closeness composition before and after the improvement, using
    the same fixed family and the two physical comparisons.
\end{proof}
